\section{Daftar: Unordered, Ordered, dan Description List}

HTML menyediakan tiga jenis elemen daftar yang masing-masing dirancang untuk keperluan spesifik dalam menyajikan informasi terstruktur. \textbf{Unordered list} (daftar tak berurutan) menggunakan elemen \code{<ul>} sebagai container dan \code{<li>} (list item) untuk setiap item. Secara default, browser menampilkan unordered list dengan bullet points (lingkaran hitam), meskipun tanda ini dapat diubah dengan CSS. Unordered list cocok untuk item yang tidak memerlukan urutan spesifik seperti daftar fitur produk, menu navigasi, atau poin-poin bahasan \cite{mdnHtmlIntro}.

\textbf{Ordered list} (daftar berurutan) menggunakan elemen \code{<ol>} sebagai container dengan \code{<li>} untuk item-itemnya. Browser menampilkan ordered list dengan nomor urut secara default. Elemen \code{<ol>} memiliki atribut berguna untuk mengontrol penomoran: atribut \code{start} menentukan angka awal penomoran (berguna untuk melanjutkan daftar dari bagian sebelumnya), atribut \code{reversed} membalik urutan penomoran dari nilai tertinggi ke terendah, dan atribut \code{type} mengubah jenis penanda (\code{1} untuk angka, \code{A} untuk huruf kapital, \code{a} untuk huruf kecil, \code{I} untuk angka Romawi kapital, \code{i} untuk angka Romawi kecil). Ordered list ideal untuk langkah-langkah prosedur, peringkat, atau informasi yang memerlukan urutan kronologis \cite{webdevHtml}.

\textbf{Description list} (atau definition list dalam HTML4) menggunakan kombinasi tiga elemen: \code{<dl>} (description list) sebagai container, \code{<dt>} (description term) untuk istilah yang didefinisikan, dan \code{<dd>} (description details) untuk deskripsi atau definisi istilah tersebut. Description list sangat berguna untuk glosarium, daftar metadata (misalnya pasangan label-nilai), atau definisi konsep. Setiap \code{<dt>} dapat diikuti oleh satu atau lebih \code{<dd>}, memungkinkan satu istilah memiliki beberapa definisi atau satu definisi menjelaskan beberapa istilah terkait.

Kemampuan menyusun daftar bersarang (nested lists) memungkinkan pembuatan outline hierarki yang kompleks. Daftar bersarang dibuat dengan menempatkan elemen \code{<ul>} atau \code{<ol>} baru di dalam elemen \code{<li>}. Tingkat penyusunan bersarang teoretis tidak terbatas, namun praktiknya lebih dari tiga level mungkin membuat struktur sulit dibaca. List bersarang umum digunakan untuk menu multi-level, outline dokumen, atau struktur organisasi. Browser biasanya mengubah bullet style secara otomatis untuk level berbeda (misalnya lingkaran penuh, lingkaran kosong, kotak) untuk menciptakan hierarki visual.

\begin{table}[htbp]
\centering
\begin{tabular}{|P{3cm}|P{2.5cm}|P{6cm}|}
\hline
\textbf{Atribut} & \textbf{Nilai} & \textbf{Efek} \\
\hline
\code{type} & \code{1}, \code{A}, \code{a}, \code{I}, \code{i} & Jenis penanda: angka, huruf, atau Romawi \\
\hline
\code{start} & Angka & Angka awal penomoran \\
\hline
\code{reversed} & - (boolean) & Urutan penomoran terbalik \\
\hline
\end{tabular}
\caption{Atribut Elemen Ordered List (\code{<ol>})}
\label{tab:ol-attributes}
\end{table}

\begin{webcode}[language=HTML, caption={Contoh Berbagai Jenis Daftar}]
<!-- Unordered List -->
<h3>Fitur Produk:</h3>
<ul>
  <li>Desain ergonomis</li>
  <li>Baterai tahan lama</li>
  <li>Garansi 2 tahun</li>
</ul>

<!-- Ordered List dengan atribut -->
<h3>Langkah Instalasi:</h3>
<ol type="1" start="1">
  <li>Unduh file installer</li>
  <li>Jalankan setup.exe</li>
  <li>Ikuti wizard instalasi
    <ol type="a">
      <li>Pilih bahasa</li>
      <li>Setujui lisensi</li>
      <li>Pilih lokasi instalasi</li>
    </ol>
  </li>
  <li>Restart komputer</li>
</ol>

<!-- Description List -->
<h3>Istilah Teknis:</h3>
<dl>
  <dt>HTML</dt>
  <dd>HyperText Markup Language - bahasa markup 
      untuk membuat halaman web.</dd>
  
  <dt>CSS</dt>
  <dd>Cascading Style Sheets - bahasa untuk 
      mengatur tampilan visual halaman web.</dd>
  
  <dt>JavaScript</dt>
  <dd>Bahasa pemrograman untuk membuat halaman 
      web interaktif.</dd>
</dl>
\end{webcode}

\begin{contoh}
\textbf{Studi Kasus Mini: Outline Dokumen dengan Nested List}

Nested list ideal untuk membuat outline dokumen yang terstruktur. Berikut contoh outline untuk bab buku:

\begin{verbatim}
Bab 4: HTML Dasar
  4.1 Pengenalan HTML
    - Sejarah HTML
    - Peran di ekosistem web
  4.2 Sintaks Dasar
    - Tag dan elemen
    - Atribut
  4.3 Struktur Dokumen
    - Doctype
    - Head dan Body
\end{verbatim}

Penggunaan \code{<ol>} untuk bab dan subbab serta \code{<ul>} untuk poin detail menciptakan hierarki visual yang jelas.
\end{contoh}

